% Esquema de ensayo del puente de Palma del Río (17-07-2020): puntos de medida sobre el tablero.
\documentclass[tikz,border=4pt]{standalone}
\usepackage{resonancia}
\begin{document}
\begin{tikzpicture}[x=0.14cm, y=0.14cm]
  % tablero en planta: 0 a 65 m
  \fill[rGrisClaro] (0,0) rectangle (72,2.5);
  \fill[rGrisClaro] (0,-9.5) rectangle (72,-12);
  \draw[grueso, color=rGris] (0,2.5) -- (72,2.5);
  \draw[grueso, color=rGris] (0,-12) -- (72,-12);
  \draw[fino, color=rGris] (0,0) -- (72,0);
  \draw[fino, color=rGris] (0,-9.5) -- (72,-9.5);
  \draw[punteada] (0,-4.75) -- (72,-4.75);
  % interrupción del dibujo (el tablero sigue)
  \draw[fino, color=rGris, decorate, decoration={zigzag, segment length=3pt, amplitude=1pt}] (72,2.5) -- (72,-12);
  % eje x = 65 m
  \draw[punteada] (65,5) -- (65,-14.5);
  % puntos 1-5
  \foreach \x/\p in {0/1, 16.25/2, 32.5/3, 48.75/4, 65/5} {
    \fill[rVerde] (\x,0) circle (0.9);
    \draw[rVerde, line width=0.5pt] (\x,0) circle (1.6);
    \node[etiqueta, above=3pt] at (\x,2.5) {\p};
    \draw[fino] (\x,-12) -- (\x,-15);
  }
  % cotas
  \foreach \a/\b in {0/16.25, 16.25/32.5, 32.5/48.75, 48.75/65} {
    \draw[<->, fino] (\a,-14.5) -- (\b,-14.5);
  }
  \node[etiqueta, font=\sffamily\scriptsize] at (8.125,-17.2) {16,25 m};
  \node[etiqueta, font=\sffamily\scriptsize] at (24.4,-17.2) {16,25 m};
  \node[etiqueta, font=\sffamily\scriptsize] at (40.6,-17.2) {16,25 m};
  \node[etiqueta, font=\sffamily\scriptsize] at (56.9,-17.2) {16,25 m};
  \node[etiqueta, anchor=east, color=rGris, align=right] at (63.5,-7.2) {\scriptsize eje\\\scriptsize $x=65$ m};
  % direcciones de medida de la tabla del ensayo
  \draw[flecha, color=rNaranja] (86,0.5) -- node[above, font=\sffamily\scriptsize] {longitudinal} (104,0.5);
  \draw[flecha, color=rNaranja] (86,-4.5) -- node[right, font=\sffamily\scriptsize] {transversal} (86,-11);
\end{tikzpicture}
\end{document}
